\section{Computational verification and reproducible artifacts}
\label{verification:section}

The package separates exact finite certificates from numerical
diagnostics. Its proofs apply to unbounded derivative orders and
denominators. The programs verify finite presentations, constants,
signs, conventions, and selected asymptotic predictions; none uses
unsuccessful integer-relation searches as a proof of independence.
All reported runs completed successfully.

\subsection{Exact arithmetic checks}

\begin{itemize}
\item \path{code/verify_ranks.py} checks $486$ exact rational
presentations. The $54$ denominators are $1,\ldots,48$ and
$60,64,72,90,105,120$. The power weights use
$s=-4,-2,-1,0,1,2,3,5$; a ninth weight system includes zero and
non-power values. The checks compare prime and all-divisor row
spans, endpoint conventions, both reflection signs, and the
nonresonant unit-symbol basis.
\item \path{code/verify_jet_ranks.py} checks $160$ expanded finite-jet
presentations, including nilpotent weights. Exact fraction elimination
confirms the dimensions before and after all endpoint layers are
prescribed. The rational specializations test the algebraic theorem;
they do not approximate logarithms as rational numbers and then claim
to compute an arithmetic rank.
\item \path{code/verify_shuffle_audit.py} constructs the $13$ by $10$
shuffle-product matrix, verifies its rank $7$, and verifies the three
integer annihilators in Section~\ref{audit:poly:shuffle-rank}. It checks
that every individual coordinate is outside that row space.
\item \path{code/verify_cubic_class_numbers.py} uses only integer
arithmetic to check all finite premises of
Proposition~\ref{corr:cubic-class-numbers}: irreducibility,
discriminants, maximality tests, splitting types, and generator norms.
The deduction of class number one uses the stated Dedekind and
Minkowski theorems.
\end{itemize}

The JSON files in \path{data/} retain the exact parameters, rational
matrix information, and witnesses. In particular, the class-number
calculation includes degree-two prime ideals within the Minkowski
bound; checking only rational primes with degree-one generators
without accounting for their partners would leave a gap.

\subsection{Independent numerical representations}

The Stieltjes script uses $85$ decimal digits. It compares the
Laplace kernel with scaled Hurwitz-jet evaluation in $8$ cases,
computes $40$ moving zeros for $n=1,2,3,4$ and
$k=8,32,128,512$, and verifies the global first-index bracket at
$k=1,2,4,8,16,32,64,128$. The largest kernel discrepancy, divided
by the larger of $1$ and the comparison magnitude, was less than
$1.7\cdot10^{-86}$. This is a residual at the selected working
precision, not a certified bound on the error in every printed digit.

There is a practical numerical issue at large $k$. The raw value of
$\zeta(k+1,a)$ can be extremely small on the moving-zero scale
$a\asymp k$. An absolute-accuracy zeta routine may therefore return
an adequate unscaled answer that loses many relative digits when
multiplied by $a^{k+1}$. The supplied program avoids that loss by
scaling \emph{inside} Euler--Maclaurin summation. It evaluates
\[
 \frac{a^{k+1}}{k!}F_{n,k}(a)
 =n![z^n]\left\{
     \frac{(1+z)_k}{k!}\,a^{k+1}\zeta(k+1+z,a)
                  \right\}
\]
with truncated power series in $z$ and terms of the form
$\exp(-(k+1)\log(1+m/a))$. This avoids finite differencing in the
jet variable. Each computed root is checked again after adding
$31$ terms to the finite Euler--Maclaurin sum. The stopping rule is
numerical, and explicit interval tail enclosures remain a proposed
extension.

For $n=1$, Table~\ref{verification:zeros} illustrates the improvement
from the explicit $1/(k+1)$ coefficient in
Section~\ref{zero:asymptotic-section}. Let $A_k$ denote the
leading term plus constant offset, and put $K=k+1$.
Here the next coefficient is
$c=0.0288707254210824205588\ldots$.

\begin{table}[htbp]
\centering
\small
\begin{tabular}{r r r r}
\toprule
$k$ & computed $a_k$ & $|a_k-A_k|$ & $|a_k-A_k-c/K|$\\
\midrule
$32$ & $56.558158770512722$ & $9.2190\cdot10^{-4}$ & $4.7031\cdot10^{-5}$\\
$128$ & $227.540415772647917$ & $2.2678\cdot10^{-4}$ & $2.9726\cdot10^{-6}$\\
$512$ & $911.472053968737555$ & $5.6465\cdot10^{-5}$ & $1.8632\cdot10^{-7}$\\
\bottomrule
\end{tabular}
\caption{Numerical zero locations and errors of two proved asymptotic
approximations. The displayed digits are rounded diagnostics.}
\label{verification:zeros}
\end{table}

The bridge script checks $12$ cases at $100$ decimal digits, including
nonreal primitive characters. The largest absolute residual was
less than $2.9\cdot10^{-101}$. Deliberately omitting the factor
$1/2$ or the required conjugation produces detectable errors, so
these checks test precisely the delicate parts of the formula.

The Gaussian script checks $p=1,3,5,7,9$ in three ways: accelerated
series summation, quadrature, and the finite reduction.
At $90$ digits, their largest pairwise discrepancy was less than
$4\cdot10^{-91}$. It also checks five real or complex generating
function arguments and the resummation at $z=1$.
The correction scripts supply $11$ checks at $70$ digits for the
log-gamma and Herglotz identities, and $7$ checks at $80$ digits
for the modified-polylogarithm and root-of-unity witnesses.
The latter reproduce both the corrected and printed values of the
complex-embedding $P_4$ example.

\subsection{Build and integration}

The principal deliverables are
\path{polylogarithms_research.tex} and
\path{polylogarithms_research.pdf}. The main \LaTeX{} file is
assembled from the editable fragments in \path{sections/}; it
contains the full text and bibliography and depends only on the
included figure and ordinary \LaTeX{} packages.
\path{BUILD.sh} rebuilds it using \texttt{pdflatex}.
\path{code/run_checks.py} reruns the verification programs.
The exact Python package versions are recorded in
\path{requirements.txt} and \path{ENVIRONMENT.json}.

The correction register records proposed changes to existing files;
the package does not silently replace those source drafts.
\path{PROVENANCE.json} fixes the inspected repository state, and
\path{SHA256SUMS} records checksums of the delivered artifacts.
The figures and data are reproducible from the included programs.
This makes each proof, correction, and finite witness independently
reviewable during repository integration.
